\section*{अध्याय \ref{ch:g11:drift-to-spreading} --- महाद्वीपीय विस्थापन से समुद्री तल प्रसरण तक}

\begin{solution}{exo:g11:drift-to-spreading:1}
\omterm{def:g11:drift-to-spreading:spreading}{समुद्री तल प्रसरण} किसी मध्य-महासागरीय कटक के शिखर पर नई महासागरीय तली का बनना है,
जहाँ पिघली चट्टान ऊपर उठकर जम जाती है, और फिर उसका दोनों ओर कटक से
दूर चलना है। पूर्ण दर वह चाल है जिससे दोनों प्लेटें एक-दूसरे से दूर होती हैं;
अर्ध दर वह चाल है जिससे एक ओर की तली अक्ष से दूर चलती है। सममित
प्रसरण के लिए पूर्ण दर $= 2
\times$ अर्ध दर।
\end{solution}

\begin{solution}{exo:g11:drift-to-spreading:2}
अपने \omterm{def:g11:drift-to-spreading:remanence}{क्यूरी ताप}, लगभग \qty{580}{\celsius}, से ऊपर मैग्नेटाइट अपना स्वयं का
चुंबकन नहीं रख सकता; केवल उससे नीचे ठंडा होने पर ही उसके कण क्षेत्र के अनुदिश
चुंबकित होते हैं और वह दिशा बनाए रखते हैं। \omterm{def:g11:drift-to-spreading:remanence}{क्यूरी ताप} से ऊपर, \qty{700}{\celsius} तक
फिर गरम करने पर चट्टान अपना पुराना
अभिलेख खो देती है; फिर से ठंडी होने पर वह उस नए समय का क्षेत्र अंकित करती है।
\end{solution}

\begin{solution}{exo:g11:drift-to-spreading:3}
ध्वनि नीचे जाती है और लौटती है: $d = 1500 \times 5.2/2 = \qty{3900}{m}$।
\end{solution}

\begin{solution}{exo:g11:drift-to-spreading:4}
वर्तमान सामान्य अंतराल \num{0.78} मिलियन वर्ष पहले शुरू हुआ: वह
\num{780000} वर्षों से चल रहा है। हारामियो घटना $1.07 - 0.99 = \num{0.08}$
मिलियन वर्ष, यानी \num{80000} वर्ष, चली। पाँच मिलियन वर्ष से नई सीमाएँ
गिनने पर (0.78, 0.99, 1.07, 1.77, 1.95, 2.14, 2.15, 2.581, 3.04, 3.11,
3.22, 3.33, 3.58, 4.18, 4.29, 4.48, 4.62, 4.80, 4.89, 4.98): बीस उत्क्रमण।
\end{solution}

\begin{solution}{exo:g11:drift-to-spreading:5}
$65\,000/40\,000 \approx 1.6$: कटक पृथ्वी के चारों ओर लगभग डेढ़
बार लिपट जाते।
\end{solution}

\begin{solution}{exo:g11:drift-to-spreading:6}
$20/0.78 \approx \qty{25.6}{km}$ प्रति मिलियन वर्ष, यानी लगभग
\qty{2.6}{cm} प्रति वर्ष की अर्ध दर; पूर्ण दर लगभग \qty{5.1}{cm} प्रति वर्ष है।
\end{solution}

\begin{solution}{exo:g11:drift-to-spreading:7}
पश्चिम: $45/3.58 \approx 12.6$~km प्रति मिलियन वर्ष, प्रति वर्ष \qty{1.26}{cm};
पूर्व: $44/3.58 \approx 12.3$~km प्रति मिलियन वर्ष, प्रति वर्ष \qty{1.23}{cm}।
पूर्ण दर: $1.26 + 1.23 \approx \qty{2.5}{cm}$ प्रति वर्ष। दोनों अर्ध दरों में
लगभग \qty{2}{\percent} का अंतर है: प्रसरण लगभग पूरी तरह सममित है।
\end{solution}

\begin{solution}{exo:g11:drift-to-spreading:8}
प्रति वर्ष \qty{1.9}{cm} का अर्थ प्रति मिलियन वर्ष \qty{19}{km} है: $990/19 \approx 52$
मिलियन वर्ष, सबसे पुराने अवसाद के 49 मिलियन वर्षों से थोड़ा अधिक
(लगभग \qty{6}{\percent})।
\end{solution}

\begin{solution}{exo:g11:drift-to-spreading:9}
प्रति मिलियन वर्ष \qty{20}{km} पर: \qty{15.6}{km}, \qty{19.8}{km},
\qty{21.4}{km}, \qty{35.4}{km} और \qty{39.0}{km}। अक्ष से: 0 से
\qty{15.6}{km} तक सामान्य, \qty{19.8}{km} तक उत्क्रमित, \qty{21.4}{km} तक सामान्य
(हारामियो), \qty{35.4}{km} तक उत्क्रमित, \qty{39.0}{km} तक सामान्य (ओल्डुवाई),
फिर उत्क्रमित।
\end{solution}

\begin{solution}{exo:g11:drift-to-spreading:10}
$64 \times 10^6/4.6 \times 10^9 \approx \num{0.014}$, यानी पृथ्वी की आयु का लगभग
\qty{1.4}{\percent}। यह हेस के इस तर्क का समर्थन करता है कि
महासागरीय तली इसलिए नई है कि वह लगातार बनती और नष्ट होती रहती है: पृथ्वी के
इतिहास के अधिकांश भाग से कोई महासागरीय तली बची नहीं है।
\end{solution}

\begin{solution}{exo:g11:drift-to-spreading:11}
किसी भी क्षण तली अक्ष पर उस क्षण की ध्रुवता के साथ बनती है,
और बँट जाती है: आधी हर ओर एक ही चाल से दूर चलती है। इसलिए एक ही समय बनी
दो पट्टियाँ अक्ष से एक ही दूरी पर होती हैं, एक हर ओर,
एक ही ध्रुवता के साथ: हर ओर दूसरी ओर की प्रति है, अक्ष में
परावर्तित। यदि तली एक ओर दूसरी ओर की अपेक्षा अधिक तेज़ी से दूर चले
(असममित प्रसरण), या अक्ष कूदकर किसी नई
स्थिति में चला जाए, तो सममिति टूट जाती है।
\end{solution}

\begin{solution}{exo:g11:drift-to-spreading:12}
(1) \qty{5}{cm} का अर्थ $5 \times 10^{-5}$~\unit{km} है: हर वर्ष
$65\,000 \times 5 \times 10^{-5} = \qty{3.25}{km^2}$ की एक पट्टी जुड़ती है।
(2) $2.9 \times 10^8/3.25 \approx 9 \times 10^7$ वर्ष, लगभग 90 मिलियन
वर्ष। (3) यह सबसे पुरानी तली की आयु (190 मिलियन वर्ष) के समान परिमाण-क्रम का है
और उससे कम: महासागरीय तली के इतनी नई बने रहने के लिए,
कहीं न कहीं उतनी ही तली नष्ट होनी चाहिए जितनी कटकों पर बनती है।
\end{solution}

\begin{solution}{exo:g11:drift-to-spreading:13}
घटना \num{0.08} मिलियन वर्ष चली। प्रति मिलियन वर्ष \qty{12.5}{km} पर
उसकी धारी $0.08 \times 12.5 = \qty{1}{km}$ चौड़ी है; प्रति मिलियन वर्ष \qty{75}{km}
पर, $0.08 \times 75 = \qty{6}{km}$। \qty{3}{km} ऊपर, समुद्र की सतह से
देखने पर, अपनी गहराई से बहुत संकरी धारी केवल एक दुर्बल
क्षेत्र उत्पन्न करती है, जो विपरीत ध्रुवता वाली उसकी चौड़ी पड़ोसी धारियों के क्षेत्रों के साथ
धुँधला जाता है; अपनी गहराई से दुगुनी चौड़ी धारी साफ़ उभर आती है।
\end{solution}

\begin{solution}{exo:g11:drift-to-spreading:14}
स्थल 16: $176/11 \approx \qty{16}{m}$ प्रति मिलियन वर्ष; स्थल 20: $72/67
\approx \qty{1.1}{m}$ प्रति मिलियन वर्ष। अवसाद एक स्थान से दूसरे स्थान पर और एक समय से
दूसरे समय में बहुत अलग-अलग दरों से जमा होता है, इसलिए उसकी मोटाई
तली की आयु नहीं देती। आयु बेसाल्ट पर सीधे पड़े सबसे पुराने
अवसाद के जीवाश्मों से मिलती है।
\end{solution}

\begin{solution}{exo:g11:drift-to-spreading:15}
(1) पश्चिम: $58/2.581 \approx \qty{22.5}{km}$ प्रति मिलियन वर्ष,
प्रति वर्ष \qty{2.25}{cm}; पूर्व: $70/2.581 \approx \qty{27.1}{km}$ प्रति मिलियन
वर्ष, प्रति वर्ष \qty{2.71}{cm}; पूर्ण दर $128/2.581 \approx \qty{49.6}{km}$ प्रति
मिलियन वर्ष, लगभग \qty{5.0}{cm} प्रति वर्ष। (2) लगभग \qty{225}{km} पश्चिम और
\qty{271}{km} पूर्व। (3) उत्क्रमण के समय दोनों सीमाएँ अक्ष पर
एक साथ थीं; तब से वे पूर्ण दर से एक-दूसरे से दूर हुई हैं, चाहे
गति दोनों ओर कैसे भी बँटी हो। उनकी दूरी को
उनकी आयु से भाग देने पर पूर्ण दर मिलती है।
\end{solution}

\begin{solution}{pb:g11:drift-to-spreading:1}
\textbf{1.} स्थल 16 सबसे निकट है (\qty{191}{km}, 11 मिलियन वर्ष), स्थल 20
सबसे दूर (\qty{1270}{km}, 67 मिलियन वर्ष): कटक से जितनी दूर,
तली उतनी पुरानी।

\textbf{2.} बेसाल्ट बनते ही अवसाद उस पर बैठना शुरू हो जाता है,
इसलिए सबसे पुरानी परत लगभग बेसाल्ट जितनी पुरानी है। वह केवल
थोड़ी कम पुरानी हो सकती है, कभी अधिक पुरानी नहीं।

\textbf{3.} किलोमीटर प्रति मिलियन वर्ष में: स्थल 16: $191/11 \approx 17.4$;
स्थल 15: $380/24 \approx 15.8$; स्थल 18: $506/26 \approx 19.5$; स्थल 17:
$643/33 \approx 19.5$; स्थल 14: $727/40 \approx 18.2$; स्थल 19: $990/49
\approx 20.2$; स्थल 20: $1270/67 \approx 19.0$।

\textbf{4.} $129.5/7 \approx \num{18.5}$~km प्रति मिलियन वर्ष, लगभग
\qty{1.85}{cm} प्रति वर्ष।

\textbf{5.} हाँ: उनके अनुपात (19.5) पश्चिमी स्थलों के अनुपातों के बीच पड़ते हैं।
तली दोनों ओर लगभग एक ही दर से फैली है: प्रसरण
सममित है।

\textbf{6.} $t = d/u$: आयु दूरी के समानुपाती है।

\textbf{7.} $u = 208\,276/10\,952 \approx \num{19.0}$~km प्रति मिलियन वर्ष,
यानी प्रति वर्ष \qty{1.90}{cm}।

\textbf{8.} लगभग $2 \times 1.9 = \qty{3.8}{cm}$ प्रति वर्ष।

\textbf{9.} हाँ: प्रति वर्ष \qty{1.9}{cm} प्रति वर्ष \qty{2}{cm} से थोड़ा
कम है।

\textbf{10.} $1270/19.0 \approx 67$ मिलियन वर्ष, उसके सबसे पुराने
अवसाद जितनी।

\textbf{11.} $191/9 \approx \num{21.2}$~km प्रति मिलियन वर्ष,
प्रति वर्ष \qty{2.1}{cm}।

\textbf{12.} $1270/71 \approx \num{17.9}$~km प्रति मिलियन वर्ष,
प्रति वर्ष \qty{1.8}{cm}।

\textbf{13.} $180\,329/9854 \approx \num{18.3}$~km प्रति मिलियन वर्ष,
प्रति वर्ष \qty{1.83}{cm}।

\textbf{14.} $(1.90 - 1.83)/1.83 \approx \qty{4}{\percent}$: दोनों दरें
आयुओं और दूरियों की अनिश्चितताओं के भीतर सहमत हैं।

\textbf{15.} एक आयु अवसाद के जीवाश्मों से आती है, दूसरी बेसाल्ट के
चुंबकत्व और उत्क्रमणों के काल-मापक्रम से। एक विधि की कोई भूल
दूसरी में नहीं दिखेगी: उनकी सहमति \omterm{def:g11:drift-to-spreading:spreading}{समुद्री तल प्रसरण} की
एक कड़ी परीक्षा है।

\textbf{16.} $13.8 + 13.8 = \ang{27.6}$।

\textbf{17.} $111.2 \times \cos 30^\circ \approx \qty{96.3}{km}$ प्रति अंश;
$27.6 \times 96.3 \approx \qty{2660}{km}$।

\textbf{18.} $2660/19.0 \approx 140$ मिलियन वर्ष।

\textbf{19.} $2660/130 \approx \num{20.5}$~km प्रति मिलियन वर्ष, लगभग
प्रति वर्ष \qty{2.0}{cm}, प्रति वर्ष \qty{1.9}{cm} से थोड़ा अधिक। चूँकि पिछले
67 मिलियन वर्ष लगभग \qty{1.9}{cm} प्रति वर्ष की दर से चले, आरंभिक खुलना
थोड़ा तेज़ था: स्थल 20 और किनारे के बीच की तली, $2660 - 1270 =
\qty{1390}{km}$, $130 - 67 = 63$ मिलियन वर्षों में बनी, यानी लगभग
प्रति वर्ष \qty{2.2}{cm}।

\textbf{20.} \textbf{\ang{30} दक्षिण पर अफ़्रीका के पास की सबसे पुरानी महासागरीय तली
लगभग 130 मिलियन वर्ष पुरानी है} (आज की दर के बहिर्वेशन से 140 मिलियन वर्ष):
दक्षिण अटलांटिक वहाँ क्रिटेशस कल्प में खुलना शुरू हुआ।
\end{solution}
